\chapter{Variables transversales et multi-actifs}\label{ch:rs:cross-sectional-and-cross-asset-features}

Une entreprise vend une grande partie de sa production à un seul client. Ce client annonce une mauvaise nouvelle, et les perspectives de l'entreprise viennent de se dégrader
avec lui~; son action devrait baisser le jour même. Cohen et Frazzini, à partir d'une base de données sur les principaux clients des entreprises, ont constaté
que les cours des actions n'intègrent pas rapidement les nouvelles concernant des entreprises liées, ce qui engendre des mouvements de cours ultérieurs prévisibles~: une
stratégie long--short fondée sur cet effet a rapporté des alphas mensuels supérieurs à 150 points de base. Une variable construite à partir d'autres titres
pose une question que l'historique du titre lui-même ne peut pas poser~: qui bouge en premier, et l'autre suit-il~? Ce chapitre le mesure
de trois façons, à trois échelles de temps~: des portefeuilles par taille aux États-Unis sur soixante ans de données hebdomadaires~; un marché synthétique avec des liens client--fournisseur implantés
sur plusieurs mois~; et deux instruments simulés, dont l'un suit l'autre avec un retard d'une fraction de seconde implanté, à l'échelle des ticks.
Une même leçon revient à chaque échelle de temps~: le passé du meneur doit prédire le suiveur au-delà du passé propre du suiveur.

\section{Résidus par pairs et par secteur}

\begin{definition}[Groupe de pairs, rendement relatif aux pairs]\label{def:rs:cross-sectional-and-cross-asset-features:peer}
Un \emph{groupe de pairs}\index{groupe de pairs} d'un titre est un ensemble de titres choisis pour partager ses moteurs communs~: sa branche,
son secteur, ou ses plus proches voisins par corrélation des rendements ou par description de l'activité. Son \emph{rendement relatif
aux pairs}\index{rendement relatif aux pairs} est son rendement moins la moyenne des rendements des autres membres de son groupe sur la même
période.
\end{definition}

Retrancher les pairs est la résidualisation la moins coûteuse, cas particulier du \omterm{def:rs:anatomy-of-a-predictor:residual}{rendement résiduel} du chapitre~6, où la moyenne des pairs
est l'unique facteur, d'exposition égale à un. Le mot \emph{autres} compte.

\begin{proposition}[Un titre n'est pas son propre pair]\label{prop:rs:cross-sectional-and-cross-asset-features:loo}
Dans un groupe de $n$ titres, le rendement diminué de la moyenne du groupe, \emph{y compris} le titre, est égal à $(n-1)/n$ fois le rendement
diminué de la moyenne des $n-1$ autres.
\end{proposition}

\begin{proof}
Avec $m_{-i}$ la moyenne des autres, la moyenne complète vaut $(x_i + (n-1)m_{-i})/n$, et $x_i - (x_i + (n-1)m_{-i})/n =
\frac{n-1}{n}(x_i - m_{-i})$.
\end{proof}

L'atténuation est sans conséquence pour un classement mais non pour un niveau~: dans un groupe de quatre, un \omterm{def:rs:cross-sectional-and-cross-asset-features:peer}{rendement relatif aux pairs} qui inclut le
titre est trop faible d'un quart, et l'erreur varie selon la taille des groupes, si bien qu'un seuil ou une régression poolée traite les groupes inégalement.
\texttt{firm.leadlag.peer\_relative} exclut le titre.

Les pairs comptent pour les variables transversales construites à partir des rendements d'\emph{autres} titres. Lorsqu'un client et son fournisseur partagent un
facteur (leur branche, le marché), la part du mouvement du client qui relève du facteur atteint le fournisseur le jour même et est intégrée aux cours~;
seule la nouvelle propre au client peut arriver en retard. Le rendement du client relatif aux pairs isole cette nouvelle. Son apport dépend de
la part commune du mouvement du client~; la section suivante en mesure un exemple.

\section{Relations d'avance--retard à horizon quotidien}

\begin{definition}[Relation d'avance--retard, lien économique]\label{def:rs:cross-sectional-and-cross-asset-features:leadlag}
Deux titres ont une \emph{relation d'avance--retard}\index{relation d'avance--retard} lorsque les rendements passés de l'un (le meneur)
prédisent les rendements futurs de l'autre (le suiveur) au-delà de ce que prédit le passé propre du suiveur. Un \emph{lien
économique}\index{lien économique} est une relation d'affaires documentée entre deux entreprises (client et fournisseur, concurrents,
partenaires, une société mère et sa filiale) qui donne une raison d'en attendre une.
\end{definition}

La définition exige une prévisibilité au-delà du passé propre du suiveur, la causalité au sens de Granger (livre~4, chapitre~20), et la
clause n'est pas une formalité. Lo et MacKinlay (1990) ont mis en évidence d'importantes \omterm{def:rs:cross-sectional-and-cross-asset-features:leadlag}{relations d'avance--retard} entre les titres américains~;
Hou (2007) les localise au sein des branches, les grandes entreprises menant les petites. La bibliothèque de données de Kenneth French contient des rendements quotidiens de
quintiles de taille américains depuis 1926, à partir desquels \texttt{rs\_fetch\_size.py} ne stocke que des statistiques dérivées (la bibliothèque porte un
avis de droits d'auteur et ne comporte aucune licence). À l'échelle hebdomadaire, la corrélation du rendement du quintile de petite taille avec celui du quintile de grande taille
de la semaine précédente valait 0.31 de 1962 à 1987, 0.23 de 1988 à 2006 et 0.06 de 2007 à mi-2026~; la corrélation inverse, le petit menant le grand,
valait 0.06, $-0.09$ et $-0.05$ (la \cref{fig:rs:cross-sectional-and-cross-asset-features:size}).

\begin{omfigure}
\begin{tikzpicture}
\begin{axis}[width=11cm,height=5.4cm,ybar,bar width=10pt,ylabel={corrélation},tick label style={font=\small},
  label style={font=\small},xtick=data,xticklabels={1962--1987,1988--2006,2007--2026},ymin=-0.15,ymax=0.45,
  enlarge x limits=0.25,grid=major,grid style={gray!25},legend columns=1,
  legend style={font=\footnotesize,at={(0.5,-0.2)},anchor=north,draw=none,fill=none,
  /tikz/every even column/.append style={column sep=8pt}},table/col sep=comma]
\addplot[fill=omDef!55,draw=omDef] table[x=k,y=small_on_big]{figdata/research/10-cross-sectional-and-cross-asset-features/size.csv};
\addplot[fill=omThm!45,draw=omThm] table[x=k,y=big_on_small]{figdata/research/10-cross-sectional-and-cross-asset-features/size.csv};
\addplot[fill=gray!40,draw=gray] table[x=k,y=own_small]{figdata/research/10-cross-sectional-and-cross-asset-features/size.csv};
\legend{petit sur la semaine précédente du grand,grand sur la semaine précédente du petit,petit sur sa propre semaine précédente}
\end{axis}
\end{tikzpicture}

\omcaption{\omterm{def:rs:cross-sectional-and-cross-asset-features:leadlag}{Relation d'avance--retard} hebdomadaire entre le quintile de plus petite taille et celui de plus grande taille des actions américaines (pondération égale), sur trois périodes~:
la corrélation du rendement de chaque quintile avec celui de l'autre, et avec son propre rendement de la semaine précédente. Données issues de la Kenneth R. French
Data Library (portefeuilles de taille quotidiens, fichier CRSP 202607).}\label{fig:rs:cross-sectional-and-cross-asset-features:size}
\end{omfigure}

Les barres grises sont le piège. Le rendement du petit quintile était aussi corrélé à sa propre semaine précédente (0.37) qu'à celui du grand
quintile (0.31), et une régression de la semaine du petit quintile sur les deux semaines précédentes donne au grand quintile une statistique $t$
valant 1.9 de 1962 à 1987, $-0.7$ de 1988 à 2006 et 1.1 depuis~: sur des semaines, l'essentiel de l'avance célèbre tient à l'autocorrélation propre des petites valeurs
du type d'autocorrélation de court terme que Boudoukh, Richardson et Whitelaw (1994) ont soutenu avoir été surestimée
dans la littérature et attribuée très probablement à des facteurs institutionnels. À l'échelle quotidienne, l'avance est plus nette~: de 1962 à 1987, le quintile de grande taille de la veille
a une statistique $t$ de 13.4 au-delà de l'autocorrélation propre du petit quintile~; depuis 2007, elle vaut $-4.4$. Dans les deux cas, elle s'est estompée.

\subsection{Liens économiques, implantés}

Un marché synthétique peut contenir des liens dont la vérité est connue. \texttt{firm.synthmkt} avec \texttt{link\_share = 0.3} désigne 30~\% des titres cotés comme fournisseurs
d'un client unique tiré au hasard, et ajoute au rendement de chaque fournisseur 10~\% des chocs spécifiques de son client,
répartis uniformément sur les 21 séances suivantes~: un retard d'attention dans l'esprit de Cohen et Frazzini. Autour des annonces
de résultats du client (la \cref{fig:rs:cross-sectional-and-cross-asset-features:event}), le client bouge de 4.5~\% par unité de surprise le
jour même, et le fournisseur dérive vers 0.39~\% sur les 21 jours suivants.

\begin{omfigure}
\begin{tikzpicture}
\begin{axis}[width=11cm,height=5.4cm,xlabel={séances à partir de l'annonce du client},
  ylabel={rendement cumulé par unité de surprise (\%)},tick label style={font=\small},label style={font=\small},xmin=-5,xmax=40,
  ymin=-0.2,ymax=0.6,grid=major,grid style={gray!25},legend pos=south east,legend style={font=\footnotesize},
  table/col sep=comma]
\addplot[omThm,thick,dashed] table[x=day,y=customer_tenth]{figdata/research/10-cross-sectional-and-cross-asset-features/event.csv};
\addplot[omDef,very thick] table[x=day,y=supplier]{figdata/research/10-cross-sectional-and-cross-asset-features/event.csv};
\legend{client ($\times 0.1$),fournisseur}
\end{axis}
\end{tikzpicture}

\omcaption{Rendements cumulés ajustés du marché autour des annonces de résultats des clients, exprimés en pentes sur la surprise réelle, pour
11\,100 événements clients et 9\,654 événements fournisseurs~: le mouvement du client (divisé par dix, la réponse implantée) et celui de ses fournisseurs.
Le fournisseur y parvient en un mois. Données~: \texttt{firm.synthmkt}, \texttt{link\_share = 0.3}, graine 1.}\label{fig:rs:cross-sectional-and-cross-asset-features:event}
\end{omfigure}

La variable se définit ainsi~: à chaque fin de mois, le rendement client cumulé du fournisseur sur le mois écoulé. Son IC de rang avec le mois
suivant du fournisseur vaut 0.035 sur 107 mois ($t = 5.6$), et un portefeuille long--short des quintiles supérieur et inférieur rapporte 0.90~\% par mois, avec un
ratio de Sharpe annualisé de 1.9. La vérité (la part du lien planté dans le rendement attendu de chaque fournisseur sur le mois à venir) a un
IC de 0.074~: le dernier mois du client est un proxy bruité de la part de sa nouvelle information qui reste à diffuser, et la moitié de l'information
se perd dans le proxy. Le rendement client relatif aux pairs fait à peine mieux (0.036), car le simulateur n'implante le lien que sur des chocs spécifiques
et ses facteurs sectoriels sont faibles~; dans des données où clients et fournisseurs partagent de forts mouvements communs, l'écart est
la raison de l'utiliser.

\subsection{Extraire les liens sans la liste}

Cohen et Frazzini disposaient de la liste des clients. Sans elle, un chercheur peut classer chaque paire ordonnée de titres selon la corrélation
du rendement de l'un avec le rendement précédent de l'autre, et espérer que les vrais liens remontent en tête. Sur le marché synthétique, les 698 titres
cotés pendant les dix années forment 486\,506 paires ordonnées, parmi lesquelles 133 sont des liens plantés~:

\begin{center}
\small
\begin{tabular}{@{}lccc@{}}
\toprule
liens parmi les\ldots & 100 premiers & 1\,000 premiers & 10\,000 premiers\\
\midrule
rendements quotidiens, retard d'un jour & 0 & 0 & 3\\
rendements mensuels, retard d'un mois & 0 & 2 & 8\\
au hasard (espérance) & 0.03 & 0.27 & 2.7\\
\bottomrule
\end{tabular}
\end{center}

La corrélation retardée mensuelle d'une vraie paire vaut en moyenne 0.052, alors que les corrélations de toutes les paires ont un écart type de 0.097~: un
lien se situe à un demi-écart type au-dessus du bruit, et parmi un demi-million de paires, c'est le bruit qui l'emporte. La liste des clients retrouve les 133
liens par construction. Extraire un réseau d'avance--retard relève du problème des tests multiples du chapitre~20 dans sa forme la plus pure~; le \omterm{def:rs:cross-sectional-and-cross-asset-features:leadlag}{lien économique} est
l'a priori qui rend la recherche assez restreinte pour réussir.

\section{Avance du contrat à terme sur le comptant et plateformes croisées}

\begin{definition}[Avance du contrat à terme sur le comptant, estimateur d'avance--retard]\label{def:rs:cross-sectional-and-cross-asset-features:estimator}
L'\emph{avance du contrat à terme sur le comptant}\index{avance du contrat à terme sur le comptant} est la tendance du prix d'un contrat à terme sur indice à bouger avant les prix
des composantes de l'indice et des fonds qui le répliquent. Un \emph{estimateur d'avance--retard}\index{estimateur d'avance--retard} estime le
décalage temporel $\vartheta$ auquel deux séries de prix sont le plus liées~: ici le retard qui maximise la covariance croisée de Hayashi--Yoshida
d'une série avec l'horloge de l'autre décalée de $\vartheta$, calculée à partir de toutes les observations sans les synchroniser (Hoffmann,
Rosenbaum et Yoshida, 2013).
\end{definition}

À haute fréquence, Huth et Abergel (2014) ont observé des fonctions de corrélation croisée fortement asymétriques, surtout entre contrats à terme
et actions~; les actifs les plus liquides ont tendance à mener. Ils ont atteint une précision de 60~\% dans la prévision du prochain mouvement du prix milieu de l'actif en retard à partir
du passé de l'actif en tête, nettement mieux qu'à partir du passé de l'actif en retard lui-même~; une stratégie naïve avec des ordres au marché ne pouvait pas en tirer profit
à cause de la fourchette. \texttt{firm.tape.simulate\_pair} construit le cas avec une réponse connue~: deux instruments sur un même
prix efficient, le prix du second étant retardé d'une latence plantée, chacun avec sa propre liquidité et son propre flux d'ordres. Le prix efficient
sautant environ une fois par seconde, chaque prix milieu change 0.65 fois par seconde.

\begin{omfigure}
\begin{tikzpicture}
\begin{axis}[width=11cm,height=5.2cm,xlabel={retard $\vartheta$ (secondes)},ylabel={corrélation croisée},tick label style={font=\small},
  label style={font=\small},xmin=-3,xmax=3,ymin=0.1,ymax=0.25,grid=major,grid style={gray!25},legend columns=2,
  legend style={font=\footnotesize,at={(0.5,-0.3)},anchor=north,draw=none,fill=none,
  /tikz/every even column/.append style={column sep=8pt}},table/col sep=comma]
\addplot[omDef,very thick] table[x=lag,y=fast]{figdata/research/10-cross-sectional-and-cross-asset-features/ccf.csv};
\addplot[gray,thick,dashed] table[x=lag,y=slow]{figdata/research/10-cross-sectional-and-cross-asset-features/ccf.csv};
\draw[omThm,thick] (axis cs:0.5,0.1) -- (axis cs:0.5,0.25);
\node[font=\footnotesize,omThm,anchor=west] at (axis cs:0.55,0.12) {latence plantée 0.5 s};
\legend{paire rapide (0.65 changements/s),paire lente (0.16 changements/s)}
\end{axis}
\end{tikzpicture}

\omcaption{Covariance croisée de Hayashi--Yoshida entre les prix milieux de deux instruments simulés, l'horloge du second décalée de $\vartheta$,
pour une latence plantée de 0.5 seconde~; normalisée par les variances réalisées sur une minute. La fonction de la paire rapide est plate au sommet et
centrée sur la latence~; le prix de la paire lente change trop rarement pour la montrer. Données~: \texttt{firm.tape}, une heure, graine
10.}\label{fig:rs:cross-sectional-and-cross-asset-features:ccf}
\end{omfigure}

La corrélation croisée est plate près de son sommet (la \cref{fig:rs:cross-sectional-and-cross-asset-features:ccf})~: chaque prix milieu suit son
prix efficient avec un éventail de délais, si bien que le pic est étalé sur une seconde de part et d'autre, et son argmax tombe à 0.15 seconde
pour une latence plantée de 0.5. L'étalement est symétrique lorsque les deux instruments réagissent de la même façon~; le centre de symétrie est donc une meilleure estimation que
le pic~: \texttt{lead\_symmetric} retient le retard autour duquel la fonction est la plus symétrique sur une seconde de chaque côté, et
trouve 0.5. Sur cinq latences plantées et trois graines (la \cref{fig:rs:cross-sectional-and-cross-asset-features:recovery}), le
centre se trompe de 0.053 seconde en moyenne et de 0.2 au pire, l'argmax de 0.19 en moyenne et de 0.5 au pire. Sur la paire lente,
dont les prix milieux changent toutes les six secondes, les deux échouent~: une avance plus courte que le temps entre deux changements de prix est invisible pour tout estimateur
construit à partir des changements de prix.

\begin{omfigure}
\begin{tikzpicture}
\begin{axis}[width=9cm,height=5.6cm,xlabel={latence plantée (secondes)},ylabel={avance estimée (secondes)},
  tick label style={font=\small},label style={font=\small},xmin=-0.2,xmax=2.2,ymin=-1,ymax=2.6,grid=major,grid style={gray!25},
  legend pos=north west,legend style={font=\footnotesize},table/col sep=comma]
\addplot[black!50,domain=-0.2:2.2,samples=2,forget plot] {x};
\addplot[only marks,mark=*,omDef,mark size=2.2pt] table[x=planted,y=fast]{figdata/research/10-cross-sectional-and-cross-asset-features/recovery.csv};
\addplot[only marks,mark=x,gray,mark size=3pt,thick] table[x=planted,y=slow]{figdata/research/10-cross-sectional-and-cross-asset-features/recovery.csv};
\legend{paire rapide,paire lente}
\end{axis}
\end{tikzpicture}

\omcaption{Avance estimée (centre symétrique de la corrélation croisée de Hayashi--Yoshida) en fonction de la latence plantée, trois graines chacune,
avec la diagonale. Données~: \texttt{firm.tape}, une heure par exécution.}\label{fig:rs:cross-sectional-and-cross-asset-features:recovery}
\end{omfigure}

Une variable exige des rendements échantillonnés. À chaque intervalle d'échantillonnage, le tableau donne la corrélation du rendement de l'actif suiveur avec le
rendement précédent de l'actif en tête, la corrélation inverse, et celle du suiveur avec son propre rendement précédent, pour la paire rapide avec la
latence plantée de 0.5 seconde~:

\begin{center}
\small
\begin{tabular}{@{}lccccccc@{}}
\toprule
intervalle (s) & 0.1 & 0.5 & 1 & 2 & 5 & 10 & 30\\
\midrule
suiveur sur le dernier de l'actif en tête & 0.12 & 0.39 & 0.55 & 0.59 & 0.48 & 0.26 & $-0.02$\\
actif en tête sur le dernier du suiveur & 0.11 & 0.33 & 0.43 & 0.45 & 0.31 & 0.16 & $-0.04$\\
suiveur sur son propre dernier & 0.13 & 0.31 & 0.40 & 0.44 & 0.35 & 0.22 & $-0.03$\\
\bottomrule
\end{tabular}
\end{center}

Chaque ligne est élevée, et c'est le simulateur qui en est responsable~: ses prix milieu atteignent le prix efficient par paliers sur quelques secondes, si bien que tous les rendements au prix milieu
tendent dans le même sens aux intervalles courts (les rendements réels des milieux de cotation ne le font pas, du moins pas dans cette mesure). L'avance injectée est une
\emph{asymétrie}, c'est-à-dire la première ligne moins la deuxième~: 0.13 à deux secondes, 0.17 à cinq, 0.10 à dix, rien à trente secondes~;
la première ligne doit aussi dépasser la troisième, ce qu'elle fait à deux secondes d'un écart de 0.14. L'empreinte de l'avance est une bande étroite d'intervalles d'échantillonnage, soit quelques fois la latence~; un échantillonnage plus lâche l'estompe,
et un échantillonnage plus fin voit trop peu de changements de prix. Le même test qui a démantelé l'avance sur la taille hebdomadaire s'applique~: le meneur doit l'emporter
sur le passé propre du suiveur.

\section{Signaux multi-actifs}

\begin{definition}[Variable multi-actifs]\label{def:rs:cross-sectional-and-cross-asset-features:crossasset}
Une \emph{variable multi-actifs}\index{variable multi-actifs} prédit un titre à partir des cours d'une autre classe d'actifs ou d'un autre marché~:
un contrat à terme sur indice pour ses composants, une devise pour ses exportateurs, une matière première pour ses producteurs, les taux pour les actions, la clôture d'un marché
pour l'ouverture d'un autre.
\end{definition}

Chacune a besoin d'un mécanisme, et c'est le mécanisme qui fixe l'horloge. Un arbitrage (un contrat à terme et son panier, un fonds et ses titres, la
fair value du livre~1, chapitre~21) maintient les prix ensemble et laisse une avance de quelques millisecondes à quelques secondes, disputée par les plus rapides (livre~11). Une exposition économique (le chiffre d'affaires d'un producteur en matière première, celui d'un exportateur en devise) laisse place à
une diffusion lente sur des jours ou des mois, l'échelle de temps du lien client. Deux pièges leur sont propres. Le premier tient aux horloges~: un
marché qui clôture plus tard porte une information que le marché plus précoce ne pouvait pas avoir eue, si bien qu'une clôture de Tokyo régressée sur une clôture de New York de
la même date paraît prédictive et constitue une fuite~; chaque valeur doit être horodatée à l'instant où elle est devenue connue (chapitre~3). Le second est
le passé propre du suiveur, la leçon des deux sections précédentes.

\section{Fiches de prédicteur}

\begin{predictorcard}{Momentum client}\label{pred:rs:cross-sectional-and-cross-asset-features:customer}
\sfield{Définition} Pour chaque fournisseur, le \omterm{def:rs:cross-sectional-and-cross-asset-features:peer}{rendement relatif aux pairs} de son principal client sur le dernier mois~; classement entre fournisseurs.
\sfield{Données et horodatage} Une liste client--fournisseur telle qu'elle est connue à la date (les déclarations sont déposées avec un retard~: chapitre~3)~; les
rendements des clients jusqu'à la clôture de fin de mois.
\sfield{Justification} Les investisseurs dont l'attention est limitée n'intègrent pas rapidement les nouvelles concernant les entreprises liées (Cohen et Frazzini).
\sfield{Horizon et demi-vie} Un mois~; sur le marché synthétique avec une diffusion injectée de 21 jours, IC de 0.035 (vérité~: 0.074).
\sfield{Normalisation} Classement entre fournisseurs~; neutraliser le secteur propre du fournisseur et son momentum.
\sfield{Modes de défaillance} Données de liens périmées ou regardant vers le futur~; liens qui traduisent une exposition commune aux facteurs plutôt que des nouvelles~; encombrement (crowding) une fois que
l'effet est publié (chapitre~28).
\sfield{Sources} Cohen et Frazzini (2008)~; Menzly et Ozbas (2010)~; \texttt{rs\_leadlag.customer\_momentum}.
\end{predictorcard}

\begin{predictorcard}{Dernier rendement du meneur}\label{pred:rs:cross-sectional-and-cross-asset-features:leader}
\sfield{Définition} Le rendement au prix milieu du meneur sur le dernier intervalle $\Delta$, utilisé comme prédicteur du rendement du suiveur sur l'intervalle suivant.
\sfield{Données et horodatage} Les cotations des deux instruments, horodatées sur une même horloge au moment de leur réception (la latence entre plateformes constitue le
signal et ne doit pas être pris pour un écart d'horloge).
\sfield{Justification} La formation des prix se produit d'abord là où la négociation est la moins coûteuse et la plus liquide (Huth et Abergel).
\sfield{Horizon et demi-vie} Quelques fois l'avance~; sur la paire simulée, l'excédent sur le sens inverse est maximal entre deux et
cinq secondes et disparaît à trente.
\sfield{Normalisation} Par la volatilité du suiveur sur l'intervalle.
\sfield{Modes de défaillance} L'autocorrélation propre du suiveur prise pour une avance~; une avance trop courte pour le rythme des changements de prix~; la fourchette
(le signal ne justifie peut-être pas de la franchir).
\sfield{Sources} Hoffmann, Rosenbaum et Yoshida (2013)~; Huth et Abergel (2014)~; \texttt{rs\_leadlag.interval\_table}.
\end{predictorcard}

\section{Tutoriel~: qui bouge en premier}\label{tut:rs:cross-sectional-and-cross-asset-features:tut}

\begin{tutorial}
\textbf{Objectif.} Estimer l'avance d'un instrument simulé sur un autre à partir de leurs cotations~; retrouver des liens client--fournisseur injectés
dans un marché synthétique et voir ce que la recherche de motifs trouve sans eux. \textbf{État final~:}
\cref{fig:rs:cross-sectional-and-cross-asset-features:ccf,fig:rs:cross-sectional-and-cross-asset-features:recovery,fig:rs:cross-sectional-and-cross-asset-features:event}~;
les tableaux de ce chapitre.
\begin{tutsteps}
\item \textbf{La corrélation croisée aux retards.} La somme de Hayashi--Yoshida du livre~4, avec la seconde horloge décalée, normalisée par les variances réalisées sur une minute
(les variances au niveau des ticks sont faussées par les paliers du prix).
\omcode{code/firm/leadlag/firm_leadlag.py}{31}{42}{Corrélation croisée de Hayashi--Yoshida aux retards.}\label{lst:rs:cross-sectional-and-cross-asset-features:ccf}
\item \textbf{Le centre symétrique}, parmi les retards où la fonction vaut au moins la moitié de son maximum.
\omcode{code/firm/leadlag/firm_leadlag.py}{49}{56}{Le retard de symétrie optimale.}\label{lst:rs:cross-sectional-and-cross-asset-features:centre}
\item \textbf{Pairs sans le titre lui-même} (\cref{prop:rs:cross-sectional-and-cross-asset-features:loo}).
\omcode{code/firm/leadlag/firm_leadlag.py}{71}{83}{Rendements relatifs aux pairs, sans le titre lui-même.}\label{lst:rs:cross-sectional-and-cross-asset-features:peer}
\item \textbf{Exécuter} \texttt{lead\_recovery()}, \texttt{interval\_table()}, \texttt{customer\_momentum()}, \texttt{event\_study()},
\texttt{mining()} et \texttt{fig\_leadlag.py}~; \texttt{rs\_fetch\_size.py} une fois, pour les statistiques de taille.
\end{tutsteps}
\textbf{Pour aller plus loin.} Donner aux deux instruments des vitesses de réaction différentes (un suiveur dont les teneurs de marché sont plus lents) et observer
la ccf perdre sa symétrie~; injecter des liens à l'intérieur des secteurs et comparer les rendements clients bruts et relatifs aux pairs.
\end{tutorial}

\section{Construction~: la boîte à outils avance--retard}\label{bld:rs:cross-sectional-and-cross-asset-features:leadlag}

\begin{build}
\bfield{Rôle} Mesure d'avance--retard sur données de ticks et sur panels journaliers, ainsi que la plomberie transversale (pairs, liens) dont sont bâties les
variables des livres~8 et~9 (arbitrage d'indice, paires, momentum client).
\bfield{Interface} \texttt{hy\_ccf(t1, x1, t2, x2, lags, norm\_step)}, \texttt{lead\_estimate}, \texttt{lead\_symmetric}, \texttt{lead\_lag\_ratio},
\texttt{directional\_corr}, \texttt{peer\_relative(ret, group)}, \texttt{linked\_feature(values, link)}, \texttt{lagged\_corr\_matrix},
\texttt{top\_pairs}~; \texttt{firm.synthmkt} gagne \texttt{link\_share}, \texttt{link\_beta}, \texttt{link\_days} et \texttt{Panel.customer}
(désactivé par défaut, de sorte que les résultats antérieurs ne changent pas).
\bfield{Règles} Un retard positif signifie que la seconde série suit~; les pairs n'incluent jamais le titre~; les corrélations décalées servent à classer
des candidats, jamais à les accepter.
\bfield{Tests d'acceptation} \texttt{code/firm/leadlag/tests/}~: une trajectoire brownienne observée à des instants aléatoires par deux séries à retard connu
(0, 0.7 et 2 secondes) est retrouvé à 0.21 seconde près par les deux estimateurs~; la corrélation directionnelle est unilatérale~; une fonction bruitée à sommet plat
est centrée correctement~; pairs sans le titre, vérifiés à la main~; un suiveur d'un jour injecté est la première paire trouvée par la recherche.
\bfield{Pour aller plus loin} Avance--retard selon l'heure de la journée~; une version multivariée (un meneur, plusieurs suiveurs)~; listes de liens traitées comme \omterm{def:rs:point-in-time-data-and-the-biases:lookahead}{données point-in-time}
(\texttt{firm.pit}).
\end{build}

\begin{omsources}
\item L.~Cohen et A.~Frazzini, ``Economic links and predictable returns'', \emph{Journal of Finance} 63(4), 2008.
\item L.~Menzly et O.~Ozbas, ``Market segmentation and cross-predictability of returns'', \emph{Journal of Finance} 65(4), 2010.
\item A.~W.~Lo et A.~C.~MacKinlay, ``When are contrarian profits due to stock market overreaction?'', \emph{Review of Financial Studies}
3(2), 1990.
\item K.~Hou, ``Industry information diffusion and the lead-lag effect in stock returns'', \emph{Review of Financial Studies} 20(4), 2007.
\item J.~Boudoukh, M.~P.~Richardson et R.~F.~Whitelaw, ``A tale of three schools: insights on autocorrelations of short-horizon stock
returns'', \emph{Review of Financial Studies} 7(3), 1994.
\item M.~Hoffmann, M.~Rosenbaum et N.~Yoshida, ``Estimation of the lead-lag parameter from non-synchronous data'', \emph{Bernoulli} 19(2),
2013~; N.~Huth et F.~Abergel, ``High frequency lead/lag relationships --- empirical facts'', \emph{Journal of Empirical Finance} 26, 2014.
\item Kenneth R. French Data Library, ``Portfolios formed on size [daily]''.
\end{omsources}

\section{Exercices}

\begin{exercise}[$\star$]\label{exo:rs:cross-sectional-and-cross-asset-features:1}
Un groupe de cinq actions a rapporté aujourd'hui 1~\%, 2~\%, 3~\%, 4~\% et 10~\%. Calculer le rendement de la cinquième action relatif à ses pairs, avec et
sans elle-même, et vérifier la \cref{prop:rs:cross-sectional-and-cross-asset-features:loo}.
\end{exercise}

\begin{exercise}[$\star$]\label{exo:rs:cross-sectional-and-cross-asset-features:2}
Parmi 486\,506 paires ordonnées dont 133 sont de vrais liens, combien de vrais liens un top 10\,000 aléatoire contient-il en moyenne~? Et un top
1\,000~?
\end{exercise}

\begin{exercise}[$\star$]\label{exo:rs:cross-sectional-and-cross-asset-features:3}
Le rendement d'un fournisseur inclut 10~\% des chocs spécifiques de son client, répartis uniformément sur les 21 jours suivants. Un client chute de 6~\%
à la suite d'une nouvelle spécifique aujourd'hui. De combien le rendement attendu du fournisseur dérive-t-il sur les 21 jours suivants, et par jour~?
\end{exercise}

\begin{exercise}[$\star\star$]\label{exo:rs:cross-sectional-and-cross-asset-features:4}
Deux prix browniens, $B(t) = A(t - L)$, sont échantillonnés sur une grille de pas $\Delta \ge L$. Montrer que la corrélation du rendement de $B$ avec
le rendement précédent de $A$ vaut $L/\Delta$, et que celle de leurs rendements sur le même intervalle vaut $1 - L/\Delta$.
\end{exercise}

\begin{exercise}[$\star\star$]\label{exo:rs:cross-sectional-and-cross-asset-features:5}
Pourquoi l'estimateur de Hayashi--Yoshida, plutôt que les rendements échantillonnés sur une grille, convient-il
à l'estimation d'une avance inférieure au temps moyen entre deux transactions~?
\end{exercise}

\begin{exercise}[$\star\star$]\label{exo:rs:cross-sectional-and-cross-asset-features:6}
Le rendement hebdomadaire du quintile des petites actions est corrélé à 0.31 avec le rendement précédent du grand quintile et à 0.37 avec son propre rendement. Expliquer, à l'aide d'une
régression, pourquoi l'avance peut n'apporter que peu, et comment le chapitre l'a testée.
\end{exercise}

\begin{exercise}[$\star\star\star$]\label{exo:rs:cross-sectional-and-cross-asset-features:7}
\emph{Programmation.} Conserver la variable de momentum client mensuelle, mais faire durer la diffusion injectée 5 jours, puis 63 jours
(\texttt{link\_days}). Calculer le rank IC de la variable et celui de la vérité dans chaque cas. Que dit le résultat du choix de l'horloge d'une
variable~?
\end{exercise}

\begin{exercise}[$\star\star\star$]\label{exo:rs:cross-sectional-and-cross-asset-features:8}
\emph{Trouver la faille.} «~Le rendement quotidien du Nikkei est fortement corrélé au rendement du S\&P 500 du même jour, donc nous négocions les actions japonaises
à la clôture de Tokyo sur le mouvement du S\&P.~»
\end{exercise}

\section{Problème~: Qui bouge en premier~?}

\begin{problem}[{Problème du week-end --- une avance injectée, retrouvée et chiffrée}]\label{pb:rs:cross-sectional-and-cross-asset-features:1}
Deux instruments simulés par \texttt{firm.tape.simulate\_pair} sur un même prix efficient, le second retardé de 0.5 seconde~; puis la
même chose à des latences de 0, 0.25, 1 et 2 secondes et pour trois graines~; puis une version lente de la paire.

\medskip\noindent\textbf{Partie I --- La corrélation croisée.}
\begin{enumerate}
\item À quelle fréquence chaque prix milieu change-t-il dans la paire rapide, et dans la lente~?
\item Pourquoi normaliser la covariance croisée de Hayashi--Yoshida par des variances réalisées à la minute et non par des variances au niveau du tick~?
\item Où se trouve l'argmax de la corrélation croisée de la paire rapide, et pourquoi n'est-il pas à 0.5 seconde~?
\item Que donne le centre symétrique, et sur quelle hypothèse repose-t-il~?
\item Quel est le ratio d'avance--retard de la paire rapide à 0.5 seconde, et à latence nulle~?
\end{enumerate}

\medskip\noindent\textbf{Partie II --- Récupération.}
\begin{enumerate}[resume]
\item Quelles sont les erreurs moyenne et maximale du centre et de l'argmax sur la paire rapide~?
\item Pourquoi les deux échouent-ils sur la paire lente~?
\item Que faudrait-il mesurer pour détecter une avance de 50 millisecondes sur une paire réelle~?
\end{enumerate}

\medskip\noindent\textbf{Partie III --- La variable.}
\begin{enumerate}[resume]
\item À un intervalle de deux secondes, donner la corrélation du suiveur avec le dernier rendement du meneur, celle dans le sens inverse, et celle du suiveur avec son propre rendement.
\item Pourquoi les trois sont-elles élevées sur ce simulateur~?
\item Quelle est l'empreinte de l'avance à 2, 5, 10 et 30 secondes~?
\item Pourquoi l'empreinte disparaît-elle à trente secondes, et pourquoi est-elle faible à 0.1 seconde~?
\item En quoi cela fait-il écho au résultat hebdomadaire par quintile de taille~?
\item La variable paierait-elle de franchir une fourchette d'un tick~? Que suggèrent les travaux de Huth et Abergel~?
\end{enumerate}

\medskip\noindent\textbf{Partie IV --- Liens quotidiens.}
\begin{enumerate}[resume]
\item Quels sont l'IC de momentum client, sa statistique $t$, et l'IC de la vérité~?
\item Pourquoi l'IC de la variable vaut-il environ la moitié de celui de la vérité~?
\item Combien de liens injectés figurent dans le top 10\,000 des paires extraites, au quotidien et au mensuel, comparé à l'espérance aléatoire~?
\item Énoncer le \emph{résultat nommé}~: l'avance estimée face à la latence injectée, et la part de la prédictibilité inter-plateformes qui subsiste à
chaque intervalle d'échantillonnage.
\item Lequel est le plus facile à trouver sans connaissance préalable, une avance d'une demi-seconde entre deux instruments sur un même prix, ou une avance
d'un mois entre un client et son fournisseur~? Pourquoi~?
\item En une phrase~: qu'est-ce qui fait d'une \omterm{def:rs:cross-sectional-and-cross-asset-features:leadlag}{relation d'avance--retard} une variable et non un artefact~?
\end{enumerate}
\end{problem}

\section{Questions d'entretien}

\begin{interviewq}[$\star$ \iqroles{chercheur}]\label{iq:rs:cross-sectional-and-cross-asset-features:1}
Comment construiriez-vous un \omterm{def:rs:cross-sectional-and-cross-asset-features:peer}{groupe de pairs} pour une action, et que peut-il mal tourner avec un \omterm{def:rs:cross-sectional-and-cross-asset-features:peer}{rendement relatif aux pairs}~?
\end{interviewq}

\begin{interviewq}[$\star\star$ \iqroles{chercheur, trader}]\label{iq:rs:cross-sectional-and-cross-asset-features:2}
Comment estimeriez-vous l'avance d'un contrat à terme sur un ETF à partir de données au tick~?
\end{interviewq}

\begin{interviewq}[$\star\star$ \iqroles{chercheur}]\label{iq:rs:cross-sectional-and-cross-asset-features:3}
Les rendements des grandes actions prédisent les rendements de la semaine suivante des petites. Est-ce un effet d'avance--retard~?
\end{interviewq}

\begin{interviewq}[$\star\star$ \iqroles{chercheur}]\label{iq:rs:cross-sectional-and-cross-asset-features:4}
Vous calculez les corrélations décalées de toutes les paires parmi 1\,000 actions et vous négociez les 100 premières. Qu'attendez-vous~?
\end{interviewq}

\begin{interviewq}[$\star\star$ \iqroles{chercheur, trader}]\label{iq:rs:cross-sectional-and-cross-asset-features:5}
Pourquoi l'action d'un fournisseur réagirait-elle lentement à la nouvelle de son client, et comment le testeriez-vous~?
\end{interviewq}

\begin{interviewq}[$\star\star\star$ \iqroles{chercheur}]\label{iq:rs:cross-sectional-and-cross-asset-features:6}
Deux prix, dont l'un est une copie retardée de l'autre, sont échantillonnés à l'intervalle $\Delta$. Quelle part de la covariance du suiveur avec le meneur
est prédictible, en fonction de $\Delta$ et du retard~?
\end{interviewq}
